\subsection{AMP databases}
The APD has tracked the natural AMP literature since 2003; the APD3
release~\cite{apd3} introduced a unified search and annotation interface, and
the current APD6 site distributes curated FASTA lists, including the 2024
natural-AMP list (3{,}306 entries) used here as external validation. DBAASP
takes a different angle: it curates peptides together with experimentally
measured activities against specific strains, making it the richer resource
for potency modeling but a harder one to mirror programmatically; during this
work its full-dataset endpoint returned empty responses to non-interactive
requests, so we cite it but do not benchmark against it
(Section 8). CAMP~\cite{camp} and its associated predictors remain widely
used baselines; DRAMP and LAMP aggregate overlapping records. These resources
differ in curation policy, redundancy, and label semantics, which is exactly
why cross-database evaluation --- training on UniProt KW-0929 and testing on
APD6 --- is more informative than a second random split of the same source.

\subsection{Machine-learning AMP recognition}
Early predictors used composition features with classical models: CAMP's
built-in SVM and random-forest classifiers~\cite{camp} and a long tail of
composition-plus-SVM tools. Deep approaches followed: AMPScanner applied
recurrent and convolutional networks to raw sequences~\cite{ampscanner}, and
subsequent work explored attention and pretrained protein language models.
Two methodological concerns recur across this literature. First, most
evaluations split sequences randomly, and AMP families are highly redundant
--- many database records are isoforms or orthologs --- which can inflate random-split metrics; a high AUROC alone does not
quantify memorization without a homology-aware control.
Second, negative sets are usually ``everything else,'' which conflates
unlabeled positives with true negatives. Our design choices respond directly:
cluster-aware splitting controls the first, length-matched reviewed negatives
acknowledge the second (without eliminating it), and external validation on
an independently curated database tests whether the learned function is
about AMPs or about the training corpus.

\subsection{AMP design}
Computational AMP design has used genetic algorithms, physicochemical rules,
and generative models including GANs, VAEs, and diffusion-style generators.
A common weakness is evaluation: generated sequences are often scored by the
same model that guided generation, an echo chamber that inflates apparent
quality. We keep the designer simple: simulated annealing on a discriminative
ensemble, with a proof of fixed-temperature detailed balance and an
explicit pairwise 3-mer similarity check against the reference corpus.
That check does not prove global novelty, and optimization on the scoring
model cannot independently validate its own designs. No design is claimed active; the section exists to
demonstrate a correct, reproducible prioritization mechanism.
